% !TeX program = xelatex
\documentclass[11pt,oneside]{article}
\usepackage{ProblemsetStyle}

\hypersetup{
    pdftitle={20260909 Problem set - Solutions},
    pdfauthor={김정우},
    pdfsubject={Science Quiz mathematics practice problem set solutions},
    pdfcreator={XeLaTeX}
}

\begin{document}

\problemsettitle{20260909 Problem set}

\begin{problemsolution}{1}
OpenAI가 2026년 9월 8일 solution을 proposed했다고 발표한 Clay Millennium Prize Problem은
\[
    \boxed{\text{Navier--Stokes existence and smoothness problem}}
\]
이다. 발표에는 AI-generated writeup과 Lean formal proof가 함께 공개되었다.
\end{problemsolution}

\begin{problemsolution}{2}
배각공식
\[
    (2\cos\theta)^2-2=2\cos(2\theta)
\]
에 의해 귀납적으로
\[
    a_n
    =
    2\cos\left(\frac{2^{n+1}\pi}{7}\right)
\]
을 얻는다. $2$의 거듭제곱을 $14$로 나눈 나머지는 주기 $3$을 가지며,
\[
    2^{2027}\equiv4\pmod{14}.
\]
따라서
\[
    \boxed{a_{2026}=2\cos\left(\frac{4\pi}{7}\right)}.
\]
\end{problemsolution}

\begin{problemsolution}{3}
$\mathbf{A}\mathbf{B}$와 $\mathbf{B}\mathbf{A}$는 $0$이 아닌 eigenvalue들을 algebraic multiplicity까지 동일하게 갖는다. 먼저
\[
    \begin{aligned}
        p_{\mathbf{A}\mathbf{B}}(\lambda)
        &=
        \det\!\begin{pmatrix}
            \lambda-2&-1\\
            -1&\lambda-2
        \end{pmatrix}\\
        &=(\lambda-2)^2-1\\
        &=(\lambda-1)(\lambda-3).
    \end{aligned}
\]
따라서 $3\times3$ 행렬 $\mathbf{B}\mathbf{A}$의 eigenvalue는 $0,1,3$이고,
\[
    \boxed{
        p_{\mathbf{B}\mathbf{A}}(\lambda)
        =
        \lambda(\lambda-1)(\lambda-3)
    }.
\]
\end{problemsolution}

\begin{problemsolution}{4}
반박된 conjecture는
\[
    \boxed{\text{Petersen Coloring Conjecture}}
\]
이다. 2026년 8월 Putman이 $112$-vertex counterexample을 제시했고, Jooken이 이어서 human-checkable proof를 공개했다.
\end{problemsolution}

\begin{problemsolution}{5}
Row player가 첫 번째 행을 probability $p$로 선택한다고 하자. Column player가 첫 번째 열을 고를 때와 두 번째 열을 고를 때의 expected payoff는 각각
\[
    1+2p,
    \qquad
    2-2p
\]
이다. 두 값을 같게 두면
\[
    1+2p=2-2p,
\]
이므로
\[
    p=\frac14,
    \qquad
    v=\frac32.
\]

이제 column player가 첫 번째 열을 probability $q$로 선택한다고 하자. Row player가 첫 번째 행과 두 번째 행을 고를 때의 expected payoff는 각각
\[
    3q,
    \qquad
    2-q
\]
이다. 두 값을 같게 두면
\[
    3q=2-q,
\]
이므로
\[
    q=\frac12.
\]
따라서
\[
    \boxed{
        p=\frac14,
        \qquad
        q=\frac12,
        \qquad
        v=\frac32
    }.
\]
\end{problemsolution}

\begin{problemsolution}{6}
Vi\`ete's formulas에 의해
\[
    \alpha+\beta+\gamma=0,
    \qquad
    \alpha\beta+\beta\gamma+\gamma\alpha=-3.
\]
또한 $1$은 주어진 다항식의 근이 아니므로 분모는 모두 $0$이 아니다. 통분하면 분자는
\[
    \begin{aligned}
        &(1-\beta)(1-\gamma)
        +(1-\gamma)(1-\alpha)
        +(1-\alpha)(1-\beta)\\
        &\qquad
        =3-2(\alpha+\beta+\gamma)
        +(\alpha\beta+\beta\gamma+\gamma\alpha)
        =0.
    \end{aligned}
\]
따라서
\[
    \boxed{0}.
\]
\end{problemsolution}

\newpage

\begin{problemsolution}{7}
Shou-Wu Zhang이 2026 International Congress of Basic Science에서 받은 mathematics medal은
\[
    \boxed{\text{Andrew Wiles Medal}}
\]
이다. 수상 citation은 arithmetic geometry에 대한 transformative contributions를 강조한다.
\end{problemsolution}

\begin{problemsolution}{8}
$S_6$의 orbit은 $6$개의 원소를 두 원소씩 세 쌍으로 나누는 모든 partition들로 이루어진다. 그 개수는
\[
    \frac{6!}{2^3\,3!}=15.
\]
Orbit--stabilizer theorem에 의해
\[
    \left|(S_6)_{\mathcal{P}}\right|
    =
    \frac{|S_6|}{|S_6\cdot\mathcal{P}|}
    =
    \frac{6!}{15}
    =48.
\]
따라서
\[
    \boxed{48}.
\]
\end{problemsolution}

\begin{problemsolution}{9}
주어진 설명에 해당하는 수학자는
\[
    \boxed{\text{Peter D. Lax}}
\]
이다. Courant Institute는 그가 2025년 5월 16일 $99$세로 별세했다고 발표했다.
\end{problemsolution}

\begin{problemsolution}{10}
주어진 식은
\[
    z^2-z+1=0
\]
과 동치이므로
\[
    z=e^{\pm i\pi/3}.
\]
따라서
\[
    z^m+z^{-m}
    =
    2\cos\left(\frac{m\pi}{3}\right).
\]
$2026\equiv4\pmod6$이므로
\[
    \boxed{
        z^{2026}+z^{-2026}
        =
        2\cos\left(\frac{4\pi}{3}\right)
        =-1
    }.
\]
\end{problemsolution}

\end{document}
